% !TeX root = ../../Book3.tex
\section{圆盘与上半平面模型}
\label{b3:sec:1303}

半平面中的度量元更简单，距离公式也能直接显示到实边界和无穷远的行为。Cayley 变换保证两模型的所有几何结论一致。

\subsection{圆盘模型}
\label{b3:subsec:130301}

圆盘模型适于以内点为观察中心，测地线为正交圆弧，伪双曲球仍是欧氏圆盘。0到半径 \(r\) 点的距离是 \(\log((1+r)/(1-r))\)，所以 \(r\uparrow1\) 时对数发散。

\ProseParagraphBreak
欧氏圆心不是内在球心的欧氏位置，除非内在球心为0。给定内在半径 \(s\)，伪双曲半径为 \(\tanh(s/2)\)，再用第6章的球公式才能画出正确的欧氏圆。

\subsection{半平面模型}
\label{b3:subsec:130302}

半平面长度元为 \(|\mathrm{d}z|/\operatorname{Im}z\)。距离满足
\[
 \cosh\mathrm d_{\mathbb H}(z,w)
 =1+\frac{|z-w|^2}{2\operatorname{Im}z\operatorname{Im}w}.
\]
由半平面伪双曲分式 \(\delta=|z-w|/|z-\overline w|\)，代入 \(\cosh(2\operatorname{arctanh}\delta)=(1+\delta^2)/(1-\delta^2)\)，并用两个分母平方之差 \(4\operatorname{Im}z\operatorname{Im}w\) 即得公式。

\ProseParagraphBreak
同一竖线上 \(\mathrm d(x+\mathrm{i}y_1,x+\mathrm{i}y_2)=|\log(y_2/y_1)|\)。竖线是测地线，其他测地线是以实轴点为圆心的半圆。水平直线通常不是测地线，即使沿它可以方便地计算一条路径长度。

\ProseParagraphBreak
上半平面中的竖直线段给极简计算。连接 \(\mathrm{i}y_1\) 与 \(\mathrm{i}y_2\) 的长度是 \(|\log(y_2/y_1)|\)，因为 \(\int\mathrm{d}y/y\) 正是长度积分。水平线 \(y=c\) 的一段 Euclid 长度为 \(b-a\)，双曲长度则为 \((b-a)/c\)，但一般不是两端之间的最短路。保持同一高度的长水平段与上方弯起的正交半圆相比，会累积更长的长度；路径长度公式与距离公式分别回答这两个问题。

\ProseParagraphBreak
在半平面中，平移 \(w\mapsto w+t\) 与正伸缩 \(w\mapsto a w\) 都保持 \(|\mathrm{d}w|/\operatorname{Im}w\)。第十一章的实 Möbius 变换由这些及反演组合，并满足同样不变性。半平面到圆盘的 Cayley 变换把无穷远送到一个圆周点，因此圆盘中看起来特殊的某个边界点，可能只是换坐标后的无穷远；双曲几何没有给它额外的内部地位。

\subsection{Cayley 变换}
\label{b3:subsec:130303}

对 \(C(z)=(z-\mathrm{i})/(z+\mathrm{i})\)，有 \(|C'|=2/|z+\mathrm{i}|^2\)，以及 \(1-|C|^2=4\operatorname{Im}z/|z+\mathrm{i}|^2\)。所以
\[
 \frac{2|C'(z)|}{1-|C(z)|^2}=\frac1{\operatorname{Im}z}.
\]
这同时核对密度归一化、距离式和曲率。若圆盘密度误写为 \(1/(1-|z|^2)\)，半平面密度会变成 \(1/(2y)\)，曲率也变为 \(-4\)，后面的 Ahlfors--Schwarz 常数将整体错误。

\subsection{无穷远边界}
\label{b3:subsec:130304}

圆周或扩充实轴是模型的理想边界，不属于度量空间。两个不同边界点确定一条完整测地线；一个边界点还确定一族与边界相切的欧氏圆，在半平面中为水平线，它们称为等距边界曲线或 horocycle。它们与测地线正交，但自身不最短。

\ProseParagraphBreak
例如 \(\mathrm{i}y\) 当 \(y\downarrow0\) 或 \(y\uparrow\infty\) 都逃向理想边界，与 \(\mathrm{i}\) 的距离均趋于无穷。欧氏意义的“离原点很近”和“很远”在内在几何中都可表示同一种无限远离。

\ProseParagraphBreak
圆盘边界是理想边界，不是完备化所增添的 Cauchy 极限。向圆周径向走，距离为 \(\log((1+r)/(1-r))\to\infty\)，不可能是一条趋向有限距离点的 Cauchy 路径。比如 \(r_n=1-e^{-n}\)，相邻点的双曲距离趋近1；虽然 Euclid 差趋于0，它在双曲度量下不是 Cauchy 列。两种拓扑在内部相同，趋向边界的度量行为却不同。

\ProseParagraphBreak
以半平面无穷远为理想端点的等距平移轨道 \(y=c\) 称为水平等高曲线的模型；其曲线长度随水平参数线性增加。经 Möbius 变换，它们变为与圆盘边界相切的圆。这样的圆与测地线不同：测地圆弧正交于边界，等高曲线相切于边界。两者在图形上都可能是圆弧，判别应看边界接触角和它们承担的几何任务。

\begin{proposition}[平移与伸缩的位移长度]\label{b3:prop:13-deck-displacement}
在上半平面模型中，双曲变换 \(T(w)=aw\)、\(a>1\) 的最小位移为 \(\log a\)，由虚轴上的点达到。抛物变换 \(P(w)=w+1\) 的位移下确界为0，但没有点达到。
\end{proposition}
\begin{proof}
令 \(w=x+\mathrm{i}y\)。距离公式给
\[
 \cosh d(w,aw)=1+\frac{(a-1)^2(x^2+y^2)}{2ay^2}
 \ge\frac{a+a^{-1}}2=\cosh\log a.
\]
等号恰在 \(x=0\)。平移情形则 \(\cosh d(w,w+1)=1+1/(2y^2)\)，对任何有限 \(y\) 大于1，随 \(y\to\infty\) 趋于1。
\end{proof}
循环双曲商中，虚轴下降为一条核心闭测地线，长度 \(\log a\)；第十二章的圆环模量为 \(\pi/\log a\)，与正文核心公式一致。循环抛物商则是去心盘，绕行圈可以任意短但没有最短非平凡圈。这一差别来自甲板变换的位移是否达到正最小值，而不是来自两曲面基本群不同。
